\section{Random Walk}

\begin{definition}[Random/Stochastic Process]
    \label{def:random-stochastic-process}%
    A \emph{random/stochastic process} is a sequence of random variables \(\para{X_n}_{n \in \NN}\).
\end{definition}

\begin{definition}[Random Walk]
    \label{def:random-walk}%
    A \emph{random walk} is a random process that can be expressed in the following way:
    \[
        X_n = x + \sum_{i = 1}^{n} Y_i,
    \]
    where \(x\) is a deterministic number, and \(\para{Y_i}_{n \in \NN}\)s are independent and identically distributed (i.i.d.). The random variables \(Y_i\)s the \emph{steps} of the random walk.
\end{definition}

\subsection{Simple Random Walk}

\begin{definition}[Simple Random Walk]
    \label{def:simple-random-walk}%
    A random walk is \emph{simple} if the steps can only take values \(\pm 1\).
\end{definition}

In this case, we may set \(\Prob\para{Y_i = 1} = p\) and \(\Prob\para{Y_i = -1} = q\), with \(p + q = 1\) and \(p, q \geq 0\).

\begin{notation}
    We denote \(\Prob_x\para{A} = \Prob\para{\Conditional{A}{X_0 = x}}\).
\end{notation}

\begin{example}[Random Walk as Gambler]
    \label{ex:random-walk-gambler}%
    We may think of \(\para{X_n}\) as the fortune of a gambler, who starts with \(x\) in their pocket at time \(0\), and every time step they bet \(1\), which he loses with probability \(\para{1 - p}\), and gets it back doubled with probability \(p\). The game ends if they reach \(0\) (which means they go bankrupt) or they reach \(a\), whichever happens first.

    We want to know the probability they reach \(a\) before going bankrupt. We let
    \[
        h\para{x} = \Prob_x \para{\para{X_n} \text{ reaches }a\text{ before reaching }0}.
    \]

    Note that \(h\para{0} = 0\) and \(h\para{a} = 1\). We have
    \begin{align*}
        h\para{x} &= \Prob_x \para{X \text{ reaches }a\text{ before }0} \\
        &= \Prob_x \para{\Conditional{X \text{ reaches }a\text{ before }0}{Y_1 = 1}} \cdot p + \Prob_x \para{\Conditional{X \text{ reaches }a\text{ before }0}{Y_1 = -1}} \cdot \para{1 - p} \\
        &= p \cdot h \para{x + 1} + \para{1 - p} \cdot h \para{x - 1}.
    \end{align*}
\end{example}
